\documentclass{amsart}
\usepackage{amsmath,amssymb,amsthm,hyperref}
\begin{document}

\begin{center}
\textbf{Resolution of rm-13348: no mathematical problem to prove}
\end{center}

\section*{Determination}

The source text associated with this entry (problem.tex) reads, verbatim:

\begin{quote}
``No well-defined mathematical problem can be extracted: the provided text
is a prose passage about human communication and contains no mathematical
statement, definitions, objects, or question. As such, there is no problem
to formalize.''
\end{quote}

This is not a mathematical statement, conjecture, or question. It contains:
\begin{itemize}
  \item no mathematical objects (no sets, numbers, functions, spaces,
        structures, etc.),
  \item no definitions,
  \item no hypotheses,
  \item no claim or quantity whose truth or value could be established.
\end{itemize}

\section*{Why no proof exists}

A proof is an argument establishing the truth of a precise proposition $P$.
Here there is no proposition $P$ to be established: the input is a prose
passage about human communication, and the accompanying meta-statement merely
\emph{asserts} that nothing mathematical can be formalized from it. There is
therefore no theorem, lemma, or claim for which a derivation could be
constructed.

Producing a ``proof'' of mathematical content would require fabricating a
statement that is not present in the source. That would be dishonest and
mathematically meaningless. Accordingly, the only correct and honest outcome
is to record that the task is vacuous: there is no mathematical content to
prove, disprove, or formalize.

\section*{Conclusion}

The entry rm-13348 is non-mathematical. No proof task is well posed, and none
can be carried out. This document records that determination as the final
resolution.

\end{document}
